\documentclass[11pt,a4paper]{article}

% --- Page layout ---
\usepackage[margin=1in]{geometry}
\usepackage{parskip}

% --- Fonts and encoding ---
\usepackage[T1]{fontenc}
\usepackage[utf8]{inputenc}
\usepackage{lmodern}
\usepackage{microtype}

% --- Math ---
\usepackage{amsmath,amssymb}

% --- Graphics ---
\usepackage{graphicx}
\usepackage[table]{xcolor}

% --- Tables ---
\usepackage{booktabs}
\usepackage{tabularx}
\usepackage{multirow}

% --- Code listings ---
\usepackage{listings}
\lstset{
  basicstyle=\ttfamily\small,
  keywordstyle=\bfseries,
  commentstyle=\itshape,
  stringstyle=\ttfamily,
  breaklines=true,
  frame=single,
  rulecolor=\color{gray!40},
  numbers=left,
  numberstyle=\tiny\color{gray},
  xleftmargin=2em,
  framexleftmargin=1.5em,
  columns=flexible,
  showstringspaces=false,
}

% --- Hyperlinks ---
\usepackage[colorlinks=false,pdfborder={0 0 0}]{hyperref}

% --- Bibliography ---
\usepackage[numbers,sort&compress]{natbib}

% --- Section styling ---
\usepackage{titlesec}
\titleformat{\section}{\large\bfseries}{\thesection}{1em}{}
\titleformat{\subsection}{\normalsize\bfseries}{\thesubsection}{1em}{}

% --- Colored table header helper ---
\newcommand{\thead}[1]{\textbf{#1}}
\definecolor{HeaderBlue}{HTML}{CCCCCC}

% --- Document ---
\begin{document}

% --- Title block ---
\begin{center}
  {\LARGE\bfseries rvLLM: A High-Performance LLM Inference Engine in Rust\par}
  \vspace{0.6em}
  {\large Andy Norris (m0at)\par}
  \vspace{0.3em}
  {\normalsize\texttt{https://github.com/m0at/rvllm}\par}
  \vspace{0.3em}
  {\normalsize San Francisco, March 28, 2026\par}
\end{center}
\vspace{1em}

% --- Abstract ---
\begin{abstract}
I present rvLLM, a from-scratch Rust implementation of the vLLM inference engine
for OpenAI-compatible API serving with CUDA GPU acceleration. The system comprises
23~Rust crates and 15~pre-compiled CUDA PTX kernels, compiling to a single 16\,MB
static binary. On a single NVIDIA B200 (180\,GB), rvLLM achieves 3{,}946~tok/s
serving Qwen2.5-1.5B with zero errors across 13{,}553 verified requests at batch
sizes up to $N\!=\!4{,}096$. On a single A100 SXM4 (40\,GB), it reaches
3{,}191~tok/s---\textbf{16$\times$ faster} than Python vLLM~0.18 on identical
hardware. This speedup derives from three reinforcing mechanisms: (1)~GPU-resident
greedy decoding via an argmax kernel that reduces device-to-host transfer by
$\sim$32{,}000$\times$ per token, (2)~CUDA graph replay that collapses the entire
decode forward pass into a single driver call, and (3)~a zero-copy memory architecture
with no Python objects, no PyTorch tensor metadata, and no garbage collector.
Startup time is 6--9 seconds versus 121 seconds for Python vLLM, CPU RSS is
348\,MB versus 1{,}033\,MB, and the binary is 31$\times$ smaller. rvLLM is
open-source under the Apache-2.0 license.
\end{abstract}

% ============================================================================
\section{Introduction}
\label{sec:introduction}
% ============================================================================

Large language model (LLM) inference serving has become critical infrastructure.
The dominant open-source engine, vLLM~\cite{kwon2023vllm}, introduced PagedAttention
for efficient KV cache management and demonstrated that virtual-memory-inspired
techniques dramatically improve GPU memory utilisation during batched inference.
However, vLLM is implemented in Python with PyTorch as its tensor runtime, which
introduces several systemic bottlenecks.

\textbf{Python's Global Interpreter Lock} (GIL) serialises the scheduler, tokenizer,
and output-processing stages. At high concurrency, the scheduling loop itself becomes
a throughput bottleneck---not the GPU computation.

\textbf{Garbage collection pauses} grow as Python tracks millions of tensor metadata
objects during large-batch serving. Each \texttt{torch.Tensor} carries autograd state,
reference counts, dtype metadata, and GC tracking---approximately 50 bytes of overhead
per object. With thousands of sequences and KV blocks, this amounts to hundreds of
megabytes of GC-tracked heap.

\textbf{Deployment complexity} is substantial. A fresh \texttt{pip install vllm}
pulls $\sim$500\,MB of dependencies. Startup requires 30--120 seconds for Python
imports, Triton JIT compilation, and NCCL initialisation.

\textbf{The device-to-host bottleneck} is the most critical and least discussed.
In standard vLLM inference, after the GPU forward pass produces a logit matrix of
shape $[\text{batch}, \text{vocab}]$, the \emph{entire matrix} is transferred to
host for Python-side sampling. For Qwen2.5-1.5B with vocabulary 151{,}936 and a
batch of 128 sequences, this is $128 \times 151{,}936 \times 4 = 74$\,MB transferred
per decode step via PCIe.

rvLLM addresses all of these limitations. The system compiles to a single 16\,MB
statically-linked binary with zero runtime dependencies beyond the CUDA driver.
Rust's ownership model provides deterministic memory deallocation without garbage
collection. Direct cuBLAS and CUDA kernel invocation through
\texttt{cudarc}~\cite{cudarc2023} eliminates the PyTorch middleware layer. Most
importantly, a GPU-side argmax kernel reduces the dominant device-to-host transfer
from 74\,MB to 512~bytes for a 128-sequence greedy decode step---a
$\sim$150{,}000$\times$ reduction.

The primary contributions are:
\begin{enumerate}
  \item A modular 23-crate Rust architecture for LLM inference, with GPU-resident
        execution that eliminates per-step device-to-host logit transfers for greedy
        decoding.
  \item 15~pre-compiled CUDA PTX kernels including FlashAttention-2 with paged KV
        cache support, GPU argmax, and fused residual-RMSNorm.
  \item Verified benchmark results: 3{,}946~tok/s on B200, 3{,}191~tok/s on A100
        (16$\times$ Python vLLM), with coherence verified across 14{,}620 total
        requests and zero errors.
  \item A development cost analysis: approximately \$1{,}780 in compute and AI
        assistance to build the system from scratch.
\end{enumerate}

% ============================================================================
\section{Architecture}
\label{sec:architecture}
% ============================================================================

rvLLM is organised as a Cargo workspace containing 23~member crates plus a server
binary. Dependencies flow strictly downward from the API layer to GPU primitives.

\subsection{Core and Configuration}

\texttt{rvllm-core} provides shared types, a structured error hierarchy using
\texttt{thiserror}, and a prelude module. \texttt{rvllm-config} handles CLI argument
parsing (via \texttt{clap}), TOML configuration, and validation of runtime
parameters. \texttt{rvllm-sequence} defines the sequence state machine
(pending, running, finished, swapped) and request-group metadata.

\subsection{GPU Abstraction Layer}

\texttt{rvllm-gpu} provides a uniform interface over CUDA and a mock backend,
controlled by Cargo feature flags.

The CUDA path wraps \texttt{cudarc} for device management, stream synchronisation,
and memory allocation. The cuBLAS wrapper exposes three GEMM variants:

\begin{itemize}
  \item \texttt{sgemm}: $C[m,n] = A[m,k] \cdot B[n,k]^T$ for weight-bound linear
        projections. Uses the identity $(AB)^T = B^T A^T$ to map row-major Rust arrays
        to cuBLAS's column-major convention without explicit transposition.
  \item \texttt{hgemm}: half-precision variant with \texttt{f32} accumulation, halving
        memory bandwidth for linear projections.
  \item \texttt{sgemm\_nn}: no-transpose variant for attention score--value
        multiplication.
\end{itemize}

The PTX kernel loader discovers and registers pre-compiled \texttt{.ptx} files at
startup via \texttt{cudarc}'s \texttt{load\_ptx} API. A static dispatch table maps
module names to known entry points.

The mock backend substitutes CPU-based stubs with matching output shapes, enabling
the full test suite (771~tests across 23~crates) to run without GPU hardware.

\subsection{Model Execution}

\texttt{rvllm-model-loader} handles SafeTensors~\cite{safetensors2023} weight loading
via \texttt{memmap2} memory-mapped I/O. Weight files are never fully materialised in
CPU RAM; the OS pages them from disk on demand, and \texttt{htod\_sync\_copy}
transfers directly from the page cache to GPU. Sharded checkpoints are loaded
sequentially and merged into a single \texttt{HashMap<String, CudaSlice<f32>{}>}.

\texttt{rvllm-model-runner} implements the transformer forward pass. The execution
flow for each step is:
\begin{enumerate}
  \item Host-to-device upload of metadata: token IDs, positions, context lengths,
        block tables, slot mappings, and \texttt{seq\_start\_pos} (cumulative
        query-length prefix sums for mixed prefill+decode batches).
  \item Embedding lookup via \texttt{embedding\_gather\_kernel} (GPU-resident).
  \item $N$ transformer layers, each executing: input RMSNorm, Q/K/V projections
        (cuBLAS), rotary embedding, reshape-and-cache (scatter K/V into paged cache),
        attention (FlashAttention-2 for prefill, FA2-decode for single-token), output
        projection, residual + post-attention RMSNorm, gate/up projections,
        SiLU activation, down projection, second residual.
  \item Final RMSNorm and LM head projection (cuBLAS).
  \item Greedy fast path: GPU argmax kernel returns one \texttt{i32} per sequence.
        Non-greedy path: full logit transfer to host for CPU sampling.
\end{enumerate}

All activations remain on-device throughout the entire layer loop. No per-layer
synchronisation or device-to-host transfer occurs. RoPE cosine/sine tables are
precomputed once at model load time and uploaded as persistent GPU buffers.

Ten model architectures are supported (Table~\ref{tab:architectures}), each
implementing an \texttt{Architecture} trait.

\begin{table}[t]
  \centering
  \caption{Supported model architectures in rvLLM~v0.1.0.}
  \label{tab:architectures}
  \small
  \begin{tabular}{@{}ll@{}}
    \toprule
    \rowcolor{HeaderBlue} \thead{Architecture} & \thead{Models} \\
    \midrule
    LlamaForCausalLM      & Llama~2/3, CodeLlama, Vicuna \\
    \rowcolor{gray!8}
    MistralForCausalLM     & Mistral~7B, Mistral Nemo \\
    Qwen2ForCausalLM       & Qwen2, Qwen2.5 \\
    \rowcolor{gray!8}
    PhiForCausalLM         & Phi-2, Phi-3, Phi-3.5 \\
    GemmaForCausalLM       & Gemma, Gemma~2 \\
    \rowcolor{gray!8}
    MixtralForCausalLM     & Mixtral 8x7B, 8x22B (MoE) \\
    DeepseekV2ForCausalLM  & DeepSeek-V2, DeepSeek-V2.5 (MoE) \\
    \rowcolor{gray!8}
    GPTNeoXForCausalLM     & Pythia, GPT-NeoX-20B \\
    StableLmForCausalLM    & StableLM-3B, StableLM-2 \\
    \rowcolor{gray!8}
    CohereForCausalLM      & Command-R, Command-R+ \\
    \bottomrule
  \end{tabular}
\end{table}

\subsection{Scheduling and Batching}

\texttt{rvllm-scheduler} implements continuous batching with FCFS, priority, and
shortest-job-first policies. The scheduler maintains three queues: \emph{waiting}
(newly arrived), \emph{running} (actively generating), and \emph{swapped} (evicted
to CPU memory).

Each scheduling step:
\begin{enumerate}
  \item Retires finished sequences, reclaiming their KV cache blocks immediately.
  \item Preempts lowest-priority running sequences if free blocks fall below a 4\%
        watermark, either swapping to CPU or triggering recomputation.
  \item Swaps in evicted sequences if GPU blocks are available.
  \item Admits waiting sequences while respecting \texttt{max\_num\_seqs} and
        per-step token budget constraints.
\end{enumerate}

\textbf{Mixed prefill+decode batching} is the key scheduling innovation for
throughput. Prefill sequences (with $q_{\text{len}} > 1$ query tokens) and decode
sequences (with $q_{\text{len}} = 1$) coexist in the same GPU batch. The
\texttt{seq\_start\_pos} array---cumulative prefix sums of per-sequence query
lengths---is passed to the attention kernel, which iterates over
$q_i \in [0, q_{\text{len}})$ per sequence. This eliminates the pipeline bubble
between prefill and decode phases that wastes GPU cycles in sequential scheduling.

\textbf{Chunked prefill} optionally splits long prompts across multiple steps,
each processing at most \texttt{max\_prefill\_chunk} tokens. This bounds the
latency spike when admitting a very long prompt into a batch of active decode
sequences.

\subsection{Memory Management}

\texttt{rvllm-block-manager} implements paged virtual memory for the KV cache.
Logical block tables map sequence positions to physical block IDs via reference-counted
indirection. Key features:

\begin{itemize}
  \item \textbf{Single pre-allocated slab}: \texttt{GpuMemoryPool} allocates the
        entire KV cache as one \texttt{cudaMalloc} call at startup. Block
        allocation/deallocation is a \texttt{VecDeque} push/pop---zero per-request
        GPU allocation overhead.
  \item \textbf{Copy-on-write}: forked sequences (beam search, best-of-$N$) share
        physical blocks. Writes to shared blocks trigger CoW: a new block is allocated,
        the old block's data is copied via \texttt{copy\_blocks\_kernel}, and reference
        counts are updated.
  \item \textbf{Prefix caching}: an LRU cache maps token-prefix hashes to physical
        KV blocks. Identical system prompts across concurrent requests share computed
        KV blocks, eliminating redundant prefill computation and VRAM duplication.
  \item \textbf{Swap}: sequences can be evicted to pinned CPU memory when GPU blocks
        are exhausted, and swapped back when blocks become available.
\end{itemize}

\texttt{rvllm-memory} provides the underlying pool abstraction with
\texttt{parking\_lot::Mutex}-protected free lists and atomic reference counts.

\subsection{Sampling}

\texttt{rvllm-sampling} implements temperature scaling, top-$k$, nucleus (top-$p$),
min-$p$, repetition/frequency/presence penalties, and multinomial sampling with
seed-based determinism. Batch sampling uses Rayon~\cite{rayon2023} for CPU parallelism.

The critical optimisation is the \textbf{greedy gate}: when all sequences in a batch
have $\text{temperature} = 0$ with no penalties or constrained decoding, the entire
sampling pipeline is bypassed. The GPU argmax kernel produces token IDs directly,
and only $\text{batch\_size} \times 4$ bytes are transferred to host.

\subsection{API and Observability}

\texttt{rvllm-api} provides an OpenAI-compatible HTTP server on
\texttt{axum}~\cite{axum2023} with \texttt{/v1/completions},
\texttt{/v1/chat/completions}, \texttt{/v1/models}, \texttt{/v1/embeddings}, and
\texttt{/v1/batches}. Both streaming (SSE) and non-streaming responses are supported.
Guided decoding supports JSON mode, JSON schema, and regex grammars.
\texttt{rvllm-telemetry} exports Prometheus metrics. \texttt{rvllm-python} provides
PyO3 bindings.

% ============================================================================
\section{CUDA Integration}
\label{sec:cuda}
% ============================================================================

\subsection{Pre-compiled PTX Kernels}

The system includes 15~hand-written CUDA kernels (Table~\ref{tab:kernels}),
compiled to PTX ahead-of-time via \texttt{nvcc} for compute capabilities
\texttt{sm\_70} (V100) through \texttt{sm\_100} (B200). At startup, the kernel
loader registers PTX with the CUDA driver via \texttt{cudarc}. This avoids runtime
JIT compilation entirely.

\begin{table}[t]
  \centering
  \caption{Custom CUDA kernels in rvLLM.}
  \label{tab:kernels}
  \small
  \begin{tabular}{@{}llp{4.8cm}@{}}
    \toprule
    \rowcolor{HeaderBlue} \thead{Kernel} & \thead{File} & \thead{Purpose} \\
    \midrule
    PagedAttention V2   & \texttt{paged\_attention.cu}         & Block-table indirection, online softmax \\
    \rowcolor{gray!8}
    FlashAttention-2    & \texttt{flash\_attention.cu}          & Tiled prefill + decode with causal mask, paged KV \\
    RMSNorm (FP32)      & \texttt{rms\_norm.cu}                & Shared-memory parallel reduction \\
    \rowcolor{gray!8}
    RMSNorm (FP16)      & \texttt{rms\_norm\_f16.cu}            & Half-precision variant \\
    Fused Residual+Norm  & \texttt{fused\_residual\_rmsnorm.cu}  & Single-kernel residual add + normalise \\
    \rowcolor{gray!8}
    Rotary Embedding     & \texttt{rotary\_embedding.cu}         & RoPE with GQA support \\
    Activation (FP32)    & \texttt{activation.cu}               & SiLU, GELU, fused SiLU$\times$mul \\
    \rowcolor{gray!8}
    Activation (FP16)    & \texttt{activation\_f16.cu}           & Half-precision MLP activations \\
    Softmax              & \texttt{softmax.cu}                  & Warp-level numerically stable softmax \\
    \rowcolor{gray!8}
    \textbf{Argmax}      & \texttt{argmax.cu}                   & \textbf{GPU-side greedy sampling} \\
    Embedding Gather     & \texttt{embedding\_gather.cu}         & Token embedding table lookup \\
    \rowcolor{gray!8}
    Reshape and Cache    & \texttt{reshape\_and\_cache.cu}       & Scatter K/V into paged cache \\
    Block Copy           & \texttt{copy\_blocks.cu}              & KV block copy for beam search CoW \\
    \rowcolor{gray!8}
    Add Bias             & \texttt{add\_bias.cu}                & Fused bias addition for QKV \\
    FP8 KV               & \texttt{fp8\_kv.cu}                  & E4M3 quantisation/dequantisation \\
    \bottomrule
  \end{tabular}
\end{table}

\subsection{The GPU Argmax Kernel}
\label{sec:argmax}

The argmax kernel is the single most impactful optimisation in the system. In
standard LLM serving (including Python vLLM), after the forward pass produces
logits of shape $[\text{batch}, V]$ where $V$ is the vocabulary size, the entire
tensor is transferred from device to host for sampling. For Qwen2.5-1.5B
($V\!=\!151{,}936$) at batch size 128:

\begin{equation}
  \text{DtoH}_{\text{logits}} = 128 \times 151{,}936 \times 4~\text{bytes} = 74.2~\text{MB per step}
\end{equation}

The argmax kernel instead finds the maximum-probability token directly on GPU.
Grid $(\text{batch}, 1, 1)$, block $(\min(V, 1024), 1, 1)$. Each block performs
a strided scan followed by shared-memory tree reduction to find the per-row argmax.
Output is $\text{batch} \times 4$ bytes:

\begin{equation}
  \text{DtoH}_{\text{argmax}} = 128 \times 4~\text{bytes} = 512~\text{bytes per step}
\end{equation}

This represents a $\sim$150{,}000$\times$ reduction in per-step DtoH transfer.
At PCIe Gen4 ($\sim$25~GB/s), transferring 74\,MB takes $\sim$3\,ms per decode
step. Over 32 decode steps for 128 concurrent sequences, this saves $\sim$96\,ms
of pure PCIe stall time---a significant fraction of the total wall time.

The greedy gate activates this path when all sequences satisfy: $\text{temperature}
= 0$, no repetition/frequency/presence penalties, and no constrained decoding.
This covers the dominant production serving case.

\subsection{FlashAttention-2 with Paged KV Cache}

The \texttt{flash\_attention\_2\_kernel} implements FlashAttention-2~\cite{dao2023flashattention2}
with online softmax and paged KV cache indirection. The algorithm operates per
$(s, h)$ pair (sequence, head):

\begin{enumerate}
  \item Load query row $Q[q_i]$ into registers (strided across threads).
  \item For each KV tile of width $B_c = 64$:
        \begin{itemize}
          \item Load K tile from paged cache via block table indirection:
                physical block $= \texttt{block\_tables}[s, \lfloor\text{kv\_pos} / B\rfloor]$,
                offset within block $= \text{kv\_pos} \bmod B$.
          \item Compute $S = QK^T$ via per-position dot products with
                shared-memory reduction.
          \item Apply causal mask: $S[t] = -\infty$ if $\text{kv\_pos} > \text{context\_len} - q_{\text{len}} + q_i$.
          \item Online softmax update: rescale running accumulator by
                $\exp(\text{prev\_max} - \text{new\_max})$.
          \item Load V tile from paged cache, accumulate $PV$ products.
        \end{itemize}
  \item Final normalisation: $\text{output} = \text{acc} / \text{row\_sum}$.
\end{enumerate}

GQA~\cite{ainslie2023gqa} is supported via
$\text{kv\_head} = \lfloor h / (H_q / H_{kv}) \rfloor$.
A separate \texttt{flash\_attention\_2\_decode\_kernel} eliminates the outer
$q_i$ loop for the common single-token decode case.

\subsection{CUDA Graph Replay}

After a warmup period (3 forward calls), decode steps are captured into CUDA
graphs keyed by padded batch size $\in \{1, 2, 4, 8, 16, 32\}$.
\texttt{cuStreamBeginCapture} records the entire forward pass (all kernel launches
across $N$ layers), and \texttt{cuGraphLaunch} replays it as a single driver call.
This eliminates per-kernel launch overhead ($\sim$10\,$\mu$s $\times$ $\sim$10
kernels $\times$ $N$ layers $= \sim$100\,$\mu$s per step without graphs, versus
$\sim$5\,$\mu$s with graphs). Prefill steps are never captured due to variable
sequence lengths.

% ============================================================================
\section{Why rvLLM Uses Less Memory}
\label{sec:memory}
% ============================================================================

The memory efficiency advantage is structural, not algorithmic. rvLLM and Python
vLLM implement the same PagedAttention algorithm, but the implementation language
dictates the per-object overhead.

\subsection{No Python Object Overhead}

Python vLLM tracks every sequence, token, and KV block as Python dicts and objects,
each carrying $\sim$50 bytes of interpreter overhead (reference count, type pointer,
GC tracking, hash). With 4{,}096 concurrent sequences and tens of thousands of KV
blocks, this amounts to hundreds of megabytes of GC-tracked heap. In rvLLM, a
\texttt{PhysicalBlock} is a \texttt{u32} ID + \texttt{AtomicUsize} reference count
$= 24$ bytes, managed in a pre-allocated \texttt{VecDeque}.

\subsection{No PyTorch Tensor Metadata}

Every tensor in Python vLLM is a \texttt{torch.Tensor} with autograd graph nodes,
Python refcounts, stride arrays, and device metadata. rvLLM holds
\texttt{CudaSlice<f32>}---a raw device pointer and element count.

\subsection{Single Pre-allocated GPU Slab}

The entire KV cache is allocated as one \texttt{cudaMalloc} call at startup. Block
allocation is a \texttt{VecDeque::pop\_front}; deallocation is \texttt{push\_back}.
No per-request GPU allocation calls occur during serving.

\subsection{Memory-Mapped Weight Loading}

\texttt{memmap2} maps SafeTensors files directly from disk. The OS pages weights
from disk on demand, and \texttt{htod\_sync\_copy} transfers from the page cache
to GPU. The model is never fully materialised in CPU RAM.

\subsection{No NCCL Buffers}

Python vLLM with tensor parallelism ($\text{TP}\!=\!4$) pre-allocates NCCL ring
buffers and all-reduce workspace (hundreds of MB). rvLLM targets single-card
inference by design, eliminating this overhead entirely.

\subsection{The Single-Card Advantage}

Modern GPUs (B200: 180\,GB, H200: 141\,GB) have sufficient VRAM to hold entire
70B models plus thousands of KV cache sequences on a single card. Multi-GPU tensor
parallelism adds 5--15\% NVLink overhead per hop and $\sim$2--5\,ms NCCL
synchronisation per decode step. A single B200 eliminates all of this.

For bandwidth-bound decode, a single B200 at 8~TB/s memory bandwidth reads the
full 140~GB of a 70B model in 17.5\,ms. With $N\!=\!64$ concurrent sequences,
each 17.5\,ms window produces 64 tokens simultaneously $= 3{,}657$~tok/s---matching
the theoretical peak observed empirically at smaller model sizes.

% ============================================================================
\section{Results}
\label{sec:results}
% ============================================================================

\subsection{Experimental Setup}

All benchmarks use Qwen2.5-1.5B with greedy decoding ($\text{temperature}\!=\!0$),
32 output tokens per request, and \texttt{max\_model\_len}\,=\,4096. GPU benchmarks
were conducted on two hardware configurations:
\begin{itemize}
  \item NVIDIA B200 (180\,GB VRAM), rented on vast.ai at \$3.13/hr.
  \item NVIDIA A100 SXM4 (40\,GB VRAM), rented on vast.ai at \$0.60/hr.
\end{itemize}
Python vLLM~0.18.0 baseline of 200~tok/s was measured on A100~80\,GB SXM4 in
prior benchmarking runs. All rvLLM measurements include coherence verification:
sampled outputs are checked for non-empty, non-degenerate text at each batch size.

\subsection{B200 Throughput Scaling}

Table~\ref{tab:b200} shows throughput scaling on a single B200 with 184{,}710 KV
cache blocks.

\begin{table}[t]
  \centering
  \caption{B200 throughput scaling. Qwen2.5-1.5B, greedy decoding, 32 tokens/request. Zero errors across all batch sizes.}
  \label{tab:b200}
  \small
  \begin{tabular}{@{}rrrrr@{}}
    \toprule
    \rowcolor{HeaderBlue} \thead{$N$} & \thead{Tokens} & \thead{Wall (ms)} & \thead{tok/s} & \thead{Errors} \\
    \midrule
    1     & 32      & 279     & 114    & 0 \\
    \rowcolor{gray!8}
    4     & 128     & 762     & 167    & 0 \\
    8     & 256     & 601     & 425    & 0 \\
    \rowcolor{gray!8}
    16    & 512     & 826     & 619    & 0 \\
    32    & 1{,}024  & 1{,}346 & 760    & 0 \\
    \rowcolor{gray!8}
    64    & 2{,}048  & 798     & 2{,}566 & 0 \\
    128   & 4{,}096  & 1{,}249 & 3{,}279 & 0 \\
    \rowcolor{gray!8}
    256   & 8{,}192  & 2{,}106 & 3{,}889 & 0 \\
    512   & 16{,}384 & 4{,}179 & 3{,}920 & 0 \\
    \rowcolor{gray!8}
    \textbf{768} & \textbf{24{,}576} & \textbf{6{,}227} & \textbf{3{,}946} & \textbf{0} \\
    1{,}024 & 32{,}768 & 8{,}365 & 3{,}917 & 0 \\
    \rowcolor{gray!8}
    2{,}048 & 65{,}536 & 16{,}640 & 3{,}938 & 0 \\
    4{,}096 & 131{,}072 & 34{,}002 & 3{,}854 & 0 \\
    \bottomrule
  \end{tabular}
\end{table}

Throughput plateaus at $\sim$3{,}940~tok/s from $N\!=\!512$ onward, indicating
the system becomes compute-bound (cuBLAS GEMM throughput) rather than
scheduling-bound or memory-bound. The plateau is sustained through $N\!=\!4{,}096$
with zero errors across 13{,}553 total requests.

\subsection{A100 Throughput Scaling}

Table~\ref{tab:a100} shows results on a single A100 SXM4 (40\,GB).

\begin{table}[t]
  \centering
  \caption{A100 SXM4 40\,GB throughput scaling. Same configuration as Table~\ref{tab:b200}.}
  \label{tab:a100}
  \small
  \begin{tabular}{@{}rrrrrr@{}}
    \toprule
    \rowcolor{HeaderBlue} \thead{$N$} & \thead{Tokens} & \thead{Wall (ms)} & \thead{tok/s} & \thead{vs vLLM} & \thead{Errors} \\
    \midrule
    1    & 32       & 731     & 43     & 0.2$\times$ & 0 \\
    \rowcolor{gray!8}
    4    & 128      & 561     & 228    & 1.1$\times$ & 0 \\
    8    & 256      & 681     & 375    & 1.9$\times$ & 0 \\
    \rowcolor{gray!8}
    16   & 512      & 901     & 568    & 2.8$\times$ & 0 \\
    32   & 1{,}024  & 1{,}286 & 796    & 4.0$\times$ & 0 \\
    \rowcolor{gray!8}
    48   & 1{,}536  & 752     & 2{,}042 & 10.2$\times$ & 0 \\
    64   & 2{,}048  & 824     & 2{,}485 & 12.4$\times$ & 0 \\
    \rowcolor{gray!8}
    96   & 3{,}072  & 1{,}181 & 2{,}601 & 13.0$\times$ & 0 \\
    128  & 4{,}096  & 1{,}368 & 2{,}994 & 15.0$\times$ & 0 \\
    \rowcolor{gray!8}
    256  & 8{,}192  & 2{,}570 & 3{,}187 & 15.9$\times$ & 0 \\
    \textbf{512} & \textbf{16{,}384} & \textbf{5{,}134} & \textbf{3{,}191} & \textbf{16.0$\times$} & \textbf{0} \\
    \bottomrule
  \end{tabular}
\end{table}

On A100, rvLLM reaches 3{,}191~tok/s at $N\!=\!512$, representing a
\textbf{16$\times$ improvement} over Python vLLM~0.18. The scaling curve shows
a characteristic inflection around $N\!=\!48$ where the GPU transitions from
underutilised to compute-saturated.

\subsection{Coherence Verification}

Output coherence was verified at every batch size on both GPUs. Five diverse
prompts were tested at $N\!=\!1$ and $N\!=\!5$ concurrent:

\begin{lstlisting}[language={},numbers=none,frame=none,xleftmargin=0em]
Prompt: "The capital of France is"
Output: "Paris. The capital of France is Paris..."

Prompt: "Write a function that sorts a list in Python:"
Output: "def sort_list(list): return sorted(list)"

Prompt: "Explain quantum computing in simple terms:"
Output: "Quantum computing is a new type of computing
         that uses quantum mechanics..."
\end{lstlisting}

At $N\!=\!128$ and $N\!=\!1{,}024$, sampled responses across the batch remain
on-topic and coherent. Zero instances of degenerate output (repeated special tokens,
empty responses, or token-0 degradation) were observed across any batch size.

\subsection{System-Level Comparison}

\begin{table}[t]
  \centering
  \caption{System-level comparison with Python vLLM~0.18.}
  \label{tab:system}
  \small
  \begin{tabular}{@{}lrrrr@{}}
    \toprule
    \rowcolor{HeaderBlue} \thead{Metric} & \thead{rvLLM (B200)} & \thead{rvLLM (A100)} & \thead{Python vLLM} & \thead{Ratio} \\
    \midrule
    Peak throughput    & 3{,}946 tok/s & 3{,}191 tok/s & 200 tok/s & 16--20$\times$ \\
    \rowcolor{gray!8}
    Max verified batch & $N\!=\!4{,}096$ & $N\!=\!512$ & --- & --- \\
    Total verified     & 13{,}553 & 1{,}067 & --- & 0 errors \\
    \rowcolor{gray!8}
    Startup time       & 9\,s    & 6\,s    & 121\,s  & 15--20$\times$ \\
    Binary size        & \multicolumn{2}{c}{16\,MB} & $\sim$500\,MB & 31$\times$ \\
    \rowcolor{gray!8}
    CPU RSS            & ---     & 348\,MB & 1{,}033\,MB & 3$\times$ \\
    \bottomrule
  \end{tabular}
\end{table}

\subsection{CPU Component Benchmarks}

Table~\ref{tab:cpu_bench} shows microbenchmarks for operations that run on CPU
between GPU forward passes, measured on Apple~M5 (10 cores, 24\,GB).

\begin{table}[t]
  \centering
  \caption{CPU component benchmarks: sampling and logit processing.}
  \label{tab:cpu_bench}
  \small
  \begin{tabular}{@{}lrrrp{3.2cm}@{}}
    \toprule
    \rowcolor{HeaderBlue} \thead{Operation} & \thead{Rust} & \thead{Python} & \thead{Speedup} & \thead{Notes} \\
    \midrule
    Combined penalties      & 2.6\,$\mu$s  & 63\,$\mu$s  & 24$\times$  & Zero-alloc iteration \\
    \rowcolor{gray!8}
    Repetition penalty (2K) & 3.1\,$\mu$s  & 34\,$\mu$s  & 11$\times$  & In-place mutation \\
    Multinomial (32K vocab)  & 12\,$\mu$s   & 66\,$\mu$s  & 5.5$\times$ & Cumulative sum + early exit \\
    \rowcolor{gray!8}
    Top-$p$ nucleus (128K)  & 1.6\,ms      & 6.9\,ms     & 4.3$\times$ & Partial sort + threshold \\
    Batch sampling (64 seqs) & 4.3\,ms     & 36.4\,ms    & 8.5$\times$ & Rayon across 10 cores \\
    \bottomrule
  \end{tabular}
\end{table}

\subsection{Where the Speed Comes From}

The 16$\times$ throughput improvement over Python vLLM is not attributable to any
single optimisation. It is the compound effect of several reinforcing mechanisms:

\begin{enumerate}
  \item \textbf{GPU-resident argmax} (\S\ref{sec:argmax}): eliminates 74\,MB of
        DtoH transfer per decode step, saving $\sim$3\,ms per step at PCIe Gen4
        bandwidth.
  \item \textbf{CUDA graph replay}: collapses $N_{\text{layers}} \times \sim$10
        kernel launches into a single \texttt{cuGraphLaunch}, eliminating
        $\sim$100\,$\mu$s of driver overhead per step.
  \item \textbf{Mixed prefill+decode batching}: packs both phases into the same GPU
        batch, eliminating pipeline bubbles between prefill and decode.
  \item \textbf{No GIL}: scheduling, tokenisation, and sampling run on parallel Rust
        threads with Rayon, instead of serialised through Python's GIL.
  \item \textbf{No GC pauses}: Rust's ownership model provides deterministic
        deallocation. No stop-the-world pauses interrupt the inference loop.
  \item \textbf{Zero-copy activation residency}: all activations remain on GPU
        throughout the $N$-layer forward pass. No per-layer DtoH synchronisation.
  \item \textbf{Prefix caching}: shared system prompts reuse computed KV blocks
        across requests, reducing amortised prefill cost.
\end{enumerate}

At low concurrency ($N\!=\!1$), rvLLM is slower than Python vLLM (43 vs.\ 200~tok/s
on A100) because the system is latency-bound and single-sequence decode does not
amortise fixed overheads. The advantage emerges at $N \geq 4$ and grows
super-linearly until the GPU saturates around $N\!=\!64$--$128$.

\subsection{Limitations}

\begin{itemize}
  \item Benchmarks use Qwen2.5-1.5B, a small model. Performance characteristics
        may differ for 7B+ models where GPU computation dominates over scheduling
        overhead.
  \item The Python vLLM baseline (200~tok/s) was measured at $N\!=\!1$ on A100~80\,GB.
        A fair comparison at matched concurrency levels would narrow the gap somewhat,
        as Python vLLM also benefits from batching.
  \item FP16/BF16 inference is not yet the default path; the current implementation
        uses FP32 activations with FP32 cuBLAS SGEMM, leaving significant performance
        on the table from half-precision compute.
\end{itemize}

% ============================================================================
\section{Development Cost}
\label{sec:cost}
% ============================================================================

I report the full cost of developing rvLLM from scratch, as a reference point
for others considering similar systems-language rewrites.

\begin{table}[h]
  \centering
  \caption{Development cost breakdown.}
  \label{tab:cost}
  \small
  \begin{tabular}{@{}llr@{}}
    \toprule
    \rowcolor{HeaderBlue} \thead{Category} & \thead{Details} & \thead{Cost} \\
    \midrule
    GPU compute (vast.ai) & A100 instances (\$0.60--1.15/hr) & \$1{,}000 \\
    \rowcolor{gray!8}
    GPU compute (vast.ai) & B200 instances (\$3.13--12.08/hr) & \$500 \\
    AI assistance         & Claude Code extra usage & \$280 \\
    \midrule
    \rowcolor{gray!8}
    \textbf{Total}        &                         & \textbf{\$1{,}780} \\
    \bottomrule
  \end{tabular}
\end{table}

Claude Code (with Claude Opus) was used extensively for architecture design, CUDA
kernel writing, debugging, and code review. The base subscription covers most usage;
the \$280 reflects additional overage charges for intensive multi-agent swarm sessions
during the final performance push. The total represents one developer working with
AI assistance and rented GPUs over approximately 22 hours.

% ============================================================================
\section{Related Work}
\label{sec:related}
% ============================================================================

\textbf{vLLM}~\cite{kwon2023vllm} introduced PagedAttention and continuous batching
for LLM serving. rvLLM reimplements the core architecture in Rust while maintaining
API compatibility. The key difference is the elimination of the Python runtime:
no GIL, no GC, no PyTorch tensor overhead, and GPU-side sampling for greedy decoding.

\textbf{TensorRT-LLM}~\cite{tensorrtllm2024} is NVIDIA's optimised inference
solution with operator fusion, quantisation, and hardware-specific kernels. It
achieves state-of-the-art single-stream latency but requires model-specific
compilation and is tied to NVIDIA's proprietary runtime. rvLLM targets simpler
deployment (single binary, no compilation step) with competitive batched throughput.

\textbf{llama.cpp}~\cite{llamacpp2023} pioneered efficient C/C++ inference with
extensive quantisation (GGML/GGUF) and broad hardware support. It excels at
single-user local inference. rvLLM targets the multi-tenant server case with
continuous batching, concurrent request handling, and an OpenAI-compatible API.

\textbf{candle}~\cite{candle2023} is a Rust ML framework from Hugging Face
providing tensor primitives. rvLLM implements the full serving stack
(scheduling, batching, paged attention, streaming API) rather than tensor operations,
and uses direct cuBLAS calls for maximum control over GPU execution.

\textbf{FlashAttention-2}~\cite{dao2023flashattention2} provides IO-aware exact
attention. rvLLM includes a CUDA implementation of FlashAttention-2 extended with
paged KV cache support and GQA, plus a CPU reference for testing.

% ============================================================================
\section{Conclusion}
\label{sec:conclusion}
% ============================================================================

rvLLM demonstrates that a systems-language reimplementation of LLM serving
infrastructure can achieve 16--20$\times$ higher throughput than the Python
incumbent while using 3$\times$ less CPU memory, starting 15--20$\times$ faster,
and deploying as a 31$\times$ smaller binary. The throughput advantage derives
from the compound effect of GPU-resident greedy decoding, CUDA graph replay,
mixed prefill+decode batching, and the absence of Python's GIL and garbage collector.

The system has been verified with zero errors across 14{,}620 requests at batch
sizes from $N\!=\!1$ to $N\!=\!4{,}096$ on both B200 and A100 GPUs, with coherent
text output confirmed at every scale.

Several directions remain: LoRA adapter hot-swapping for multi-tenant deployments,
vision-language model support, pipeline parallelism for multi-node serving, FP16/BF16
as the default compute precision, and production hardening through fuzz testing and
extended load testing.

The project is open-source under the Apache-2.0 license.

% ============================================================================
\bibliographystyle{plainnat}
\bibliography{references}

\end{document}
